%! root=../main.tex
\appendix %% CCS: DO NOT REMOVE

\section{Open Science} %% CCS: DO NOT REMOVE

\kjee{CCS: REMOVE the content below and replace it with appropriate content.}
\subsection{Artifacts}
\begin{itemize}
    \item Git repository including external submodule, source, build instructions, and binary hardening examples.\\
            Access: anonymized GitHub repository: \url{http://www.github.com/blablablabla}
\end{itemize}
Each submitted paper must include an Open Science appendix that:
\begin{itemize}
\item Enumerates all artifacts needed to evaluate the paper's core contributions (e.g., code, datasets, models, configuration files, scripts, documentation, benchmarks).
\item Clearly describes how the program committee can access each artifact during double-blind review (including anonymous URLs or credentials, where applicable).
% \item Explicitly justifies any artifact that cannot be shared (e.g., due to licensing restrictions, responsible disclosure concerns, safety or privacy of study subjects, or deployment risks if adversarial methods are released prematurely). When full sharing is not possible, authors are encouraged to provide partial, synthetic, or redacted artifacts that still allow reviewers to assess the methodology.
% \item In case no artifact is needed to evaluate the paper’s core contributions, the authors should state it explicitly.
\end{itemize}

\section{Ethical Considerations} %% CCS: the entire sectin MAY BE REMOVED only if it is
                                %% clear that your
                       %% paper does not raise any ethical
                       %% concerns. In case of doubts, keep the
                       %% section and explains why your paper  does
                       %% not raise any ethical concerns.

% \kjee{CCS: REMOVE the content below and replace it with appropriate content.}

The ethical ramifications of our work in binary defense is the same as for any work in binary defense.
Powerful malicious actors may use it to further their own interests, safeguarding themselves against retaliation or usurpation.
Although this is an ubiquitous risk, we believe that stronger defenses favor the forces of good more than the forces of evil.

% Authors are expected to consider the ethical implications and potential societal impact of their work. Papers that raise ethical concerns, such as those involving human subjects, user data, or real-world vulnerability analysis, must include a dedicated "Ethical Considerations" section. This section should discuss the balance of risks vs. benefits and the steps taken to minimize potential harm (e.g., responsible disclosure, data anonymization). Note that institutional (IRB/ERB) approval is neither strictly necessary nor always sufficient to demonstrate ethical conduct; we expect authors to reason about the ethics of their work beyond ensuring institutional compliance. For detailed guidance on community standards, we follow the USENIX Security'26 Ethics Policy\footnote{\url{https://www.usenix.org/conference/usenixsecurity26/call-for-papers\#ethics}}. This section does not count toward the page limit and must be placed after the 12-page main content.

\section{ARMv7 Instruction Categories}

Here we categorize all ARMv7 user-mode instructions based on their control-flow effects as described in \S~\autoref{sec:cfi_design}.
\ixb{The story here can be ``look at all these instructions whose \emph{operands} change how they need to be treated. Imagine how complex the implementation is. Imagine \emph{not} having a proof for its soundness!}

\begin{figure}[b]
    \centering
    \begin{tabular}{ll}
        Category & Instructions\\
        Unchanged block & \\
        Non-control-flow change & \\
        Direct control flow & \\
        Indirect control flow & \\
        Unsupported & \\
        \unpredictable{} & \\
        Operand specific & \\
    \end{tabular}
    \caption{Caption}
    \label{tab:placeholder}
\end{figure}

\section{Hash Function Replacement}

This appendix outlines the work needed to replace the hash function with one that decreases the space overhead seen in \autoref{fig:spec2017_size} without compromising the soundness of the rewriter.
In \autoref{sec:fm_friendly} we list FM friendly design principles.
Another design principle that helps in this case is making the desired soundness/correctness properties independent of deep properties of the implementation.
For instance, the soundness of the rewriter does not depend on the cryptographic security or collision avoidance of the hash function.